problem:  
Let \((M^n,g_0)\) be a complete noncompact Riemannian manifold with **nonnegative Ricci curvature** and **polynomial volume growth**:
\[
\operatorname{Vol}(B_r(p)) \le C(1+r)^k, \quad k \ge n.
\]
Consider a complete, properly embedded mean curvature flow \(\{\Sigma_t\}_{t\in[0,T)} \subset (M,g_0)\),
\[
\partial_t F = -H\nu,
\]
satisfying the local geometric bound
\[
|A(x,t)|^2 \le C(1+d_{g_0}(x,p))^{2\alpha}, \quad \alpha<1,
\]
and assume the flow develops a singularity at finite or infinite time \(T\): there exist points \((x_i,t_i)\) with \(t_i\nearrow T\), \(d_{g_0}(x_i,p)\to\infty\), and \(|A(x_i,t_i)|\to\infty\).

For each such sequence, let \(\lambda_i := |A(x_i,t_i)|\) and define the pointed, parabolically rescaled flows in exponential coordinates around \(x_i\):
\[
(M_i,g_i) := (M,\lambda_i^2 g_0),
\qquad
\Sigma^i_t := \exp_{x_i}^{-1}\big(\Sigma_{t_i+\lambda_i^{-2}t}\big).
\]
Assume local curvature bounds and noncollapsing conditions such that
\[
(M_i,g_i,\Sigma^i_t,x_i) \to (M_\infty,g_\infty,\Sigma_\infty(t),x_\infty)
\]
in the pointed \(C^\infty_{\mathrm{loc}}\) Cheeger–Gromov sense, where \((M_\infty,g_\infty)\) is a complete static Riemannian manifold and \(\Sigma_\infty(t)\) is an **ancient**, properly embedded mean curvature flow.

Suppose that \((M_\infty,g_\infty)\) admits a **line**, so that by the Cheeger–Gromoll splitting theorem,
\[
(M_\infty,g_\infty) \cong (\mathbb{R}\times N^{n-1}, ds^2 + h).
\]

**Conjectural Problem:**  
Conjecture that under these hypotheses, whenever the ambient limit manifold \((M_\infty,g_\infty)\) splits isometrically with an \(\mathbb{R}\)-factor, the corresponding limit flow \(\Sigma_\infty(t)\subset (M_\infty,g_\infty)\) is a **translating soliton** in the direction of the parallel vector field \(\partial_s\):
\[
H = \langle \partial_s, \nu \rangle,
\qquad
\Sigma_\infty(t) = \Phi_{t\partial_s}(\Sigma_\infty(0)).
\]
In other words, the existence of a splitting direction in the ambient limit forces the singularity model at infinity to exhibit pure translation along that direction.  
(Note: No claim is made that splitting must always occur.)

Why is it a "good" problem:  
This formulation is now logically consistent and geometrically natural. It conjectures that *whenever large-scale ambient geometry induces an \(\mathbb{R}\)-splitting in the blow-up limit, the corresponding hypersurface model must realize a translating soliton aligned with that split direction*. The problem links three central ideas:

1. **Global Riemannian structure:** the splitting phenomenon under nonnegative Ricci curvature captures the asymptotic flat direction of the manifold;  
2. **Singularity analysis in geometric PDE:** the blow-up mechanisms of mean curvature flow and the emergence of ancient solutions;  
3. **Self-similar geometry:** translators as canonical Type II singularity models.

This concise conjecture embodies the principle that *geometric noncompactness manifests analytically as translation*. Proving or disproving it would reveal how curvature conditions and asymptotic topology constrain the formation of singularities. Its scope and simplicity make it an exemplary “good” research-level problem bridging asymptotic geometry, analysis of flows, and geometric rigidity theory.